\section{Embracing Insomnia}

When I was a boy around 10 years old, my parents would drop me off at the Boston Museum of Science\footnote{\url{https://www.mos.org/}} on Saturday mornings while they went shopping. There would be a children's lecture program and afterwards I would wait for them in the library. My favorite topics were mathematics and astronomy, which were conveniently next to each other in the Dewey library classification system.

I loved to dream about far off stars, which used to have exotic names like Sirius, Antares, Betelgeuse, Canopus, Arcturus, or Regulus. Those names are now being replaced with dull standardized names. It's a moot point anyway because few stars can be recognized with all the light pollution today. The stars have gone the way of the armadillo, tortoises, racoons, etc., that used to wander around the neighborhoods.

My favorite books were the four volumes of the \emph{World of Mathematics}; the four dark blue volumes contained pages of mysterious symbols that revealed secrets of the universe to that boy. One of my favorite chapters was about Rene Descartes. I read that he was allowed to remain in bed in the morning in order to give him leisure time to meditate. It became a habit which he followed throughout his life, finding it ``conducive to intellectual profit and comfort.''

As an adult, during the cold months, he would retire to a heated room where he meditated all day. That was the source of his philosophy. Strangely, Socrates used to meditate in the snow although there are a limited number of snow days in Athens. There have been other cold weather philosophers like Nietzsche or Evola who preferred the mountains. Nevertheless, metaphysics began in warmer climates.

\paragraph{Alone with the alone}


Unlike the young Descartes, I was never allowed to tarry in bed. Hence, time for meditation had to be planned. A few years ago I developed insomnia. Not that I was a worrier who could not get to sleep, but I would wake up in the middle of the night. With nothing to do. At first, I would get up to read or watch and old movie. When sleep seemed to be returning, I would go back to bed.

Some 18 months ago, I complained to the doctor who then prescribed a strong soporific medication. After reading the side effects, I decided to keep the insomnia; after all, it isn't fatal.

That is when I learned to embrace that solitary time. The nighttime is so quiet and still. There is no sensory experience other than the flow of thoughts. I used that time to meditate on deep topics, watching the thoughts play with each other. Occasionally, fresh thoughts would be revealed from unknown sources. The benefits of such meditations exceed the reading of many books. This form of meditation is also called ``mental prayer'', as described by the \textbf{Venerable Louis of Granada} in \emph{Summa of the Christian Life}:

\begin{quotex}
Mental prayer is any form of meditation or contemplation of the tings of God ... Meditation is of inestimable benefit to the soul, for as the study of human sciences is the principal means of acquiring human wisdom, so the consideration of divine things is a very important means for attaining to divine wisdom.
\end{quotex}


Now I look forward to the time to be alone with the Alone.


\flrightit{Posted on 2023-02-05 by Cologero}

\begin{center}* * *\end{center}

\begin{footnotesize}
\begin{sffamily}
\texttt{Kaukomieli on 2023-02-06 at 01:18 said:}

We had a doctor around here who used to say to his patients that no one has ever died of insomnia but many die in their sleep.

\hfill

\texttt{Dimitri on 2023-02-06 at 13:45 said:}

\textgreater A few years ago I developed insomnia. Not that I was a worrier who could not get to sleep, but I would wake up in the middle of the night. With nothing to do. At first, I would get up to read or watch and old movie. When sleep seemed to be returning, I would go back to bed.

Funnily enough, this is actually normal, or at least it used to be. It's called ``biphasic sleep'', and it was common in medieval times. People would go to bed early in the night for their ``first sleep'', after which they'd wake up and engage in some activity for about an hour, and return to bed for their ``second sleep'' in the post-midnight hours. In the preindustrial West, most people slept in two discrete blocks.

\url{https://www.bbc.com/future/article/20220107-the-lost-medieval-habit-of-biphasic-sleep}


As it turns out this habit of yours aligns with the traditional lifestyle. Go figure.

\hfill

\texttt{Dimitri on 2023-02-06 at 14:19 said:}

I'm actually going through that article again and finding lots of fascinating details, I encourage others to read it too. It talks about how the practice was ubiquitous across all parts of the world for millennia (even during the Classical period) and still exists today in certain nonindustrial regions, and even talks about experiments where researchers actually restored people's biphasic sleep by reducing their exposure to artificial lighting before bed. Here are a few highlights that may be of interest to Gornahoor readers:

``The period of wakefulness that followed first sleep was known as ``the watch'' -- and it was a surprisingly useful window in which to get things done. ` The records describe how people did just about anything and everything after they awakened from their first sleep,' says Ekirch.\\
... \\

But the watch was also a time for religion. For Christians, there were elaborate prayers to be completed, with specific ones prescribed for this exact parcel of time. One father called it the most `profitable' hour, when -- after digesting your dinner and casting off the labours of the world -- `no one will look for you except for God'.

Those of a philosophical disposition, meanwhile, might use the watch as a peaceful moment to ruminate on life and ponder new ideas. In the late 18th Century, a London tradesman even invented a special device for remembering all your most searing nightly insights -- a `nocturnal remembrancer', which consisted of an enclosed pad of parchment with a horizontal opening that could be used as a writing guide.''

`` Historian Roger Ekirch wonders if today people might remember fewer dreams than our ancestors did, because it's less common to wake up in the middle of the night.''

\hfill

\texttt{Kaukomieli on 2023-02-09 at 12:41 said:}

Dimitri, there is actually a strength training regimen also in todays world where you sleep in two different times a day which helps to increase testosterone and growth hormone levels naturally. It is also claimed that you'll get more done and energy levels are better if one trains oneself to follow the regimen and discipline where you sleep 6 hours from 9 pm to 3 am, then again after exercise, cleaning and eating you'll take a nap for two to three hours. The training is done in 6 am, after three hours of waking and fulfilling the bodily system with right nutrients. You'll sleep again after noon right when you wish to.

I think this could very well be also be integrated into spiritual life, for the 9 to 3 sleeping is also recommended by some eastern authorities and yogis, and i believe it has been studid scientifically also that it is the best time to sleep for multiple reasons. (This is why todays 24/7 civiliation produces so much illness and disease I think; the asuric modern human is moved only by constant agitation, speed etc.)

To meditate on three in the morning right after sleeping, you'll be then awake for the last watches of the night and behold the Spiritual Sun and the morning star in the morning. Then you'll do the bodily exercise (run, ski, swim, gym, martial arts, yoga etc.) and energize yourself for the next productive day.

\hfill

\texttt{Matt on 2023-02-11 at 13:22 said:}

Reading about the longing to be ``alone with the Alone'' feels similar to the feeling I get from The Isle of the Dead; both the composition by Rachmaninoff and the painting by Arnold Böcklin that inspired it. Böcklin described the painting as, ``a dream picture: it must produce such a stillness that one would be awed by a knock on the door.''

Instead of being macabre, it makes me think of the kind of death that Tomberg is pointing to in Covenant of the Heart (aka, Lazarus Come Forth!) when he quotes Thomas speaking to the disciples, ``Let us also go, that we may die with him''. Tomberg says, ``If human life is a pendulum swing between the moment and eternity, then Lazarus had---in spirit, soul, and body---chosen eternity completely and radically. His breathing became ever more ``vertical''---that is, it took place increasingly in the direction above to below and became ever more shallow in the direction outer-inner, until it stood still, i.e., until he died.''

Echoing what was said about meditation or mental prayer, in Mysterium Coniunctionis Jung says, ``What I call coming to terms with the unconscious the alchemists called ``meditation.'' Ruland says of this: `Meditation: The name of an Internal Talk of one person with another who is invisible, as in the invocation of the Deity, or communion with one's self, or with one's good angel.\textquotesingle''

But it's not an ordinary kind of dialogue, like two people chewing the fat. It has the quality of showing up with, and listening with, one's entire being.

Jung also quotes Gerhard Dorn, who says, ``Meditative knowledge is thus the sure and undoubted resolution by expert certitude, of all manner of opinions concerning the truth... We cannot be resolved of any doubt save by experiment, and there is no better way to make it than on ourselves... We have said earlier that piety consists in knowledge of ourselves, and hence we begin to explain meditative knowledge from this also. But no man can truly know himself unless first he see and know by zealous meditation... what rather than who he is, on whom he depends, and whose he is, and to what end he was made and created, and by whom and through whom.''

Rachmaninoff, Isle of the Dead: \url{https://youtu.be/0VR4S69bdto}

\hfill

\end{sffamily}
\end{footnotesize}

